\section{AI and Conceptual Design}
\label{sec:related}

This review covers three linked needs: why AI-assisted conceptual
design is needed in aircraft and especially UAV development, which AI
tools and frameworks already exist, and why a tail-sitter VTOL UAV is a
useful validation case relative to other eVTOL architectures.

\subsection{Design need}

Aircraft conceptual design is still constrained by long, costly,
iterative workflows that must reconcile aerodynamics, structures,
propulsion, stability, manufacturing, and mission requirements early,
when most key decisions are fixed
\citep{karali2024surrogate,corrado2022,lukyanov2024}. Conventional
sequential methods depend heavily on designer experience and
low-fidelity semi-empirical models, which are especially weak for UAVs
because many UAV sizes, Reynolds-number regimes, and unconventional
layouts fall outside the historical data behind those regressions
\citep{karali2024surrogate,lukyanov2024}.

UAV design sharpens this need because UAVs can adopt tailless,
blended-wing, morphing, tandem, canard, and other nontraditional
configurations that demand broader design-space exploration and tighter
coupling between mission definition and vehicle architecture
\citep{karali2024surrogate,ivanova2025,lukyanov2024}. High-fidelity CFD
and finite-element analyses improve realism, but thousands of
evaluations across optimization loops make direct early-stage use
computationally prohibitive, motivating automation, surrogates, and
concurrent multidisciplinary design optimization (MDO) workflows
\citep{karali2024surrogate,wang2023,pliakos2025,leng2024}.

\begin{itemize}[nosep]
\item Early MDO detail is needed to reduce late development setbacks
  and resolve the conceptual-stage ``design paradox'' while design
  freedom remains high \citep{corrado2022}.
\item UAV programs need faster iteration because competition, budget
  limits, and expanding mission complexity make broad manual concept
  trades increasingly impractical
  \citep{karali2023,pliakos2025,shakeri2018}.
\item AI is increasingly relevant because modern UAVs already rely on
  autonomy, real-time decision-making, and mission adaptation under
  battery and payload constraints
  \citep{ivanova2025,sai2023,mcenroe2022}.
\end{itemize}

\subsection{AI tools and frameworks}

The literature points to a clear pattern: AI in conceptual aircraft
design is used less as a replacement for physics than as a
computational amplifier around physics-based models, MDO loops, and
design iteration. The dominant tools are surrogate models for
aerodynamics, structural mass, and sometimes radar cross-section;
evolutionary or multiobjective optimizers; automated geometry/CFD
pipelines; and data-mining or multi-fidelity schemes that reduce
coupling and computational burden
\citep{karali2024surrogate,leng2024,pliakos2025}.

Reported frameworks already support configuration selection, early
sizing, and iterative refinement. AI-supported UAV conceptual design
frameworks compare alternative configurations using learned
disciplinary models, feed optimized results back into sizing, and can
cut runtime dramatically; one deep-learning surrogate framework
reported more than three orders of magnitude computational-time
improvement over traditional methods \citep{karali2024surrogate}. The
methodological frontier is widening toward uncertainty-aware MDO,
digital twins, generative AI, and automated data-generation pipelines
that make high-quality training data easier to produce
\citep{karali2024surrogate,wang2023,du2026,pliakos2025}.

Table~\ref{tab:methodclasses} summarizes representative AI-supported
methodology classes reported for aircraft and UAV conceptual design.

\begin{table}[htbp]
\centering
\caption{Representative AI-supported methodology classes reported for
aircraft and UAV conceptual design.}
\label{tab:methodclasses}
\small
\begin{tabularx}{\textwidth}{@{}lXX@{}}
\toprule
\textbf{Method class} & \textbf{What it supports} &
\textbf{Representative evidence} \\
\midrule
ANN / deep surrogates &
Fast aerodynamics, mass, RCS prediction for concept comparison and
sizing &
AI-driven UAV frameworks use learned aerodynamics, structural mass,
and RCS models in early design \citep{karali2024surrogate} \\
\addlinespace
Evolutionary / GA / DE / multiobjective optimization &
Configuration search, wing optimization, takeoff-weight reduction,
trade-space exploration &
GA{+}ANN wing optimization and DE-based UAV optimization are
established \citep{boutemedjet2019,komarov2024} \\
\addlinespace
Multi-fidelity / data mining &
Variable screening, disciplinary decoupling, cheap-to-expensive
optimization progression &
Multi-fidelity data mining uses low-cost data first, then CFD
fine-tuning \citep{leng2024} \\
\addlinespace
Automated CFD pipelines &
Rapid geometry, meshing, solver setup, dataset generation &
Python-based preprocessing reduces turnaround to minutes with
sub-5\% deviation from expert setups \citep{pliakos2025} \\
\addlinespace
Generative / emerging AI &
Parameterization, predictive modeling, training facilitation,
constraint handling &
GenAI applications now span VAE, GAN, diffusion, and transformer
approaches \citep{du2026} \\
\addlinespace
LLM / agentic coding assistants &
Building and maintaining the analysis toolchain itself;
multidisciplinary change propagation; documentation &
This work; general-purpose agentic tools \citep{anthropic2025claudecode} \\
\bottomrule
\end{tabularx}
\end{table}

\begin{itemize}[nosep]
\item Reinforcement learning is widely used in AI-integrated UAV
  systems, while federated and distributed learning frameworks are
  emerging quickly \citep{sai2023}.
\item AI toolchains increasingly include libraries, frameworks, and
  edge/embedded deployment pathways, not just offline design
  optimizers \citep{sai2023,mcenroe2022}.
\item Current gaps remain important: many studies still focus on
  airfoils or wings rather than full-aircraft models, and robust
  treatment of uncertainty, unconventional configurations, and broader
  objective sets is still developing
  \citep{karali2024surrogate,wang2023,lukyanov2024}.
\end{itemize}

Notably, none of the surveyed frameworks employs AI at the
\emph{workflow} level --- as an agent that authors, verifies, and
evolves the design toolchain and its documentation under human
direction. That is the niche addressed by the present methodology.

\subsection{Tail-sitters versus eVTOLs}

Tail-sitter VTOL UAVs are a strong validation case because they force a
conceptual design method to address the hardest coupling in VTOL
design: hover, wing-borne cruise, and the transition between them.
Recent VTOL design studies argue that transition and climb are often
oversimplified in conceptual design even though they strongly drive
motor and battery sizing \citep{kaneko2024}. For tail-sitters
specifically, nonlinear dynamics during transition are a central
limitation, with widely varying speed and attitude, unstable
long-period behavior in some regimes, and controller sensitivity to
uncertainty \citep{yang2024}.

Against that difficulty, tail-sitters offer a simpler airframe than
many other eVTOL classes. Compared with tilt-rotor or dual-system
eVTOLs, tail-sitter eVTOLs remove tilting mechanisms and redundant
propulsors, which can reduce structural and mechanical complexity
\citep{yang2024}. But that simplicity trades for harder transition
physics: the aircraft must sustain pitch authority through a full
$90^{\circ}$ rotation, propulsors must provide both high vertical
thrust and efficient forward-flight operation, and oblique inflow
during transition can excite propeller vibration \citep{yang2024}.

\begin{itemize}[nosep]
\item Trajectory-coupled optimization materially matters for
  tail-sitters: simultaneous design-and-trajectory optimization reduced
  energy by 29\% versus uncoupled optimization in one tail-sitter
  delivery case \citep{kaneko2024}.
\item Optimized transition control improves operations, with reduced
  transition time and altitude excursion relative to traditional
  linear transition methods \citep{yang2024}.
\item A tail-sitter validation example is methodologically valuable
  because it stress-tests whether an AI-assisted conceptual design
  framework can integrate physics, controls, propulsion, and sizing in
  one loop rather than treating configuration choice and transition
  performance separately \citep{kaneko2024}.
\end{itemize}

Taken together, the literature supports an AI-assisted conceptual
design methodology that augments physics-based analysis rather than
replacing it, is especially justified for UAVs with unconventional
configurations, and is well validated by a tail-sitter VTOL case
because tail-sitters expose exactly the multidisciplinary and
transition-coupled trade-offs that simpler conceptual methods tend to
miss.
